% !TEX root = ../../tree_notes.tex
% =====================================================================
%  1.1  What Is a Tree?                                 sec:c01-what-is-tree
%  Exports: def:c01-tree  def:c01-rooted  def:c01-degree
%           thm:c01-edges-nodes  thm:c01-unique-path
%           lst:c01-build  prob:c01-same-tree
% =====================================================================
\section{What Is a Tree?}\label{sec:c01-what-is-tree}

\begin{topicintro}
A tree is a set of nodes joined by edges with no loops and no separate islands --
you can get from any node to any other, and there is only one way to do it. File
systems, organisation charts, HTML documents and parse trees are all trees, and every
data structure in this book is built on this one picture.

\textbf{Prerequisites.} The idea of a graph (nodes and edges) from a first
algorithms course.

\textbf{After this section you can}
\begin{itemize}[nosep]
  \item tell a tree from a graph that is not a tree, at a glance;
  \item root a tree and name parents, children, leaves and subtrees correctly;
  \item prove that a tree on $n$ nodes has $n-1$ edges and one path per pair;
  \item read a tree from an edge list into an adjacency list in \cpp.
\end{itemize}
\end{topicintro}

\subsection{Intuition}

\begin{analogybox}[Folders and family]
A tree is a computer's folder structure: one root folder, sub-folders inside folders,
files at the bottom, and no folder ever nesting back inside its own descendant. A family
tree drawn down from a single ancestor is the same picture. That "nothing is its own
ancestor" rule is exactly the no-cycles condition we make precise below.
\end{analogybox}

Look at the picture below. The dots are nodes and the lines are edges. There are no
closed loops: starting anywhere, you can never come back to where you began without
retracing a step. That loop-free, fully-connected shape \emph{is} a tree.

\begin{figure}[H]
\centering
\begin{tikzpicture}[diagram]
  \node[troot] (a) at (0,0) {A};
  \node[tnode] (b) at (-1.6,-1.3) {B};
  \node[tnode] (c) at (1.6,-1.3) {C};
  \node[tleaf] (d) at (-2.4,-2.6) {D};
  \node[tleaf] (e) at (-0.8,-2.6) {E};
  \node[tleaf] (f) at (1.6,-2.6) {F};
  \draw[tedge] (a) -- (b); \draw[tedge] (a) -- (c);
  \draw[tedge] (b) -- (d); \draw[tedge] (b) -- (e);
  \draw[tedge] (c) -- (f);
\end{tikzpicture}
\caption{A tree. The orange node is chosen as the \emph{root}; the green nodes
$D, E, F$ are \emph{leaves}. Notice there is exactly one path between any two
letters.}\label{fig:c01-first-tree}
\end{figure}

Now the three pictures that are \textbf{not} trees. A loop (a cycle) is forbidden, and
so is a floating piece you cannot reach.

\begin{figure}[H]
\centering
\begin{tikzpicture}[diagram]
  % cycle
  \node[numnode] (x1) at (0,0) {};
  \node[numnode] (x2) at (1,0) {};
  \node[numnode] (x3) at (0.5,-0.9) {};
  \draw[tedge] (x1) -- (x2) -- (x3) -- (x1);
  \node[annot] at (0.5,0.8) {has a cycle};
  % disconnected
  \node[numnode] (y1) at (3,0) {};
  \node[numnode] (y2) at (4,0) {};
  \node[numnode] (y3) at (3.5,-1.3) {};
  \draw[tedge] (y1) -- (y2);
  \node[annot] at (3.5,-1.9) {two pieces};
  % a real tree for contrast
  \node[numnode] (z1) at (6,0) {};
  \node[numnode] (z2) at (5.4,-1) {};
  \node[numnode] (z3) at (6.6,-1) {};
  \draw[tedge] (z1) -- (z2); \draw[tedge] (z1) -- (z3);
  \node[annot] at (6,0.8) {a tree};
\end{tikzpicture}
\caption{Left: a cycle, so not a tree. Middle: disconnected (node isolated), so not a
tree. Right: connected and loop-free, so a tree.}\label{fig:c01-not-trees}
\end{figure}

The same unrooted tree becomes a \emph{rooted} tree the moment we pick a starting node
and let gravity decide who is above whom.

\begin{figure}[H]
\centering
\begin{tikzpicture}[diagram]
  \begin{scope}[xshift=-1.6cm]
    \node[troot] (r) at (0,0) {A};
    \node[tnode] (b) at (-0.9,-1.2) {B};
    \node[tnode] (c) at (0.9,-1.2) {C};
    \node[tleaf] (d) at (-0.9,-2.4) {D};
    \draw[tedge] (r) -- (b); \draw[tedge] (r) -- (c); \draw[tedge] (b) -- (d);
    \node[annot] at (0,0.9) {rooted at A};
  \end{scope}
  \begin{scope}[xshift=2.1cm]
    \node[tnode] (r2) at (0,-0.2) {A};
    \node[troot] (b2) at (-0.9,-1.4) {B};
    \node[tnode] (c2) at (0.9,0.4) {C};
    \node[tleaf] (d2) at (-0.9,-2.6) {D};
    \draw[tedge] (r2) -- (b2); \draw[tedge] (r2) -- (c2); \draw[tedge] (b2) -- (d2);
    \node[annot] at (0.1,0.9) {rooted at B};
  \end{scope}
\end{tikzpicture}
\caption{One tree, two roots. Rooting at $B$ makes $A$ a child of $B$ and $D$ a child
of $B$; rooting at $A$ makes $B$ a child of $A$. The edges do not change, only our
view of parent and child does.}\label{fig:c01-rooting}
\end{figure}

\subsection{Definitions and Notation}

\begin{definition}{Tree and forest}{c01-tree}
A \term{tree} is a connected graph with no cycles. A \term{forest} is a disjoint union
of trees -- it drops only the ``connected'' requirement. Equivalently, a graph on $n$
nodes is a tree if any (and hence all) of these hold: it is connected with $n-1$
edges; it has no cycle but adding any missing edge creates one; every pair of nodes is
joined by exactly one path.
\end{definition}

\begin{definition}{Rooted tree: parent, child, leaf}{c01-rooted}
A \term{rooted tree} is a tree with one distinguished node called the \term{root}.
Directing every edge away from the root gives the relations: if $u$ sits above $v$
and they are adjacent, $u$ is the \term{parent} of $v$ and $v$ is a \term{child} of
$u$; a node with no child is a \term{leaf}; a node that is not a leaf is
\term{internal}. The \term{subtree} at $v$, written $\subtree(v)$, is $v$ together
with all its descendants. An \term{ancestor} of $v$ is any node on the path from the
root to $v$ (including $v$ itself); a \term{descendant} runs the other way. Siblings
share a parent.
\end{definition}

\begin{figure}[H]
\centering
\begin{tikzpicture}[diagram]
  \node[troot] (a) at (0,0) {A};
  \node[tnode] (b) at (-1.5,-1.3) {B};
  \node[tnode] (c) at (1.5,-1.3) {C};
  \node[tnode] (d) at (-2.3,-2.6) {D};
  \node[tleaf] (e) at (-0.7,-2.6) {E};
  \node[tleaf] (f) at (1.5,-2.6) {F};
  \draw[tedge] (a) -- (b); \draw[tedge] (a) -- (c);
  \draw[tedge] (b) -- (d); \draw[tedge] (b) -- (e); \draw[tedge] (c) -- (f);
  \begin{scope}[on background layer]
    \node[subtreebox, fit=(d)(b)(e)] {};
  \end{scope}
  \node[annot] at (-2.6,-1.6) {parent B};
  \node[annot] at (2.6,-1.9) {child C};
  \node[annot, text=softpurple] at (-1.5,-3.2) {subtree(B)};
\end{tikzpicture}
\caption{Anatomy of the rooted tree. $B$ is the parent of $D$ and $E$ (its children);
$C$ is a child of $A$; $D,E,F$ are leaves; the dashed box is the subtree rooted at
$B$. Node $A$ is an ancestor of every other node.}\label{fig:c01-anatomy}
\end{figure}

\begin{definition}{Degree, path, distance}{c01-degree}
The \term{degree} $\degop(v)$ of a node is its number of neighbours (in a rooted tree,
often split into in-degree $=$ children and out-degree $=$ $1$ for the parent). A
\term{path} is a sequence of nodes each adjacent to the next; its \term{length} is the
number of edges. The \term{distance} $\dist(u,v)$ is the length of the (unique) shortest
path between $u$ and $v$.
\end{definition}

\subsection{Counting Edges, Finding Paths}

Two facts about trees come up in almost every proof later in the book. They look
obvious, so it is worth seeing exactly why they hold.

\begin{theorem}{Edges of a Tree}{c01-edges-nodes}
A tree with $n$ nodes has exactly $n-1$ edges.
\end{theorem}
\begin{proofbox}
Induction on $n$. For $n=1$ there are no edges, and $1-1=0$. For $n\ge 2$, a finite
tree always has a leaf: walk a longest simple path, and one of its endpoints must have
degree $1$ -- if it had another neighbour we could either extend the path or close a
cycle, both impossible. Call that leaf $v$; deleting $v$ and its single incident edge
leaves a tree on $n-1$ nodes, which by the induction hypothesis has $n-2$ edges. Putting
$v$ back adds exactly one edge, giving $(n-2)+1=n-1$.
\end{proofbox}
\begin{insightbox}
Read it backwards: $n-1$ edges is the \emph{minimum} needed to connect $n$ nodes. One
edge short and you have two islands; one edge too many and you create a cycle. That is
why contest statements say ``$n$ cities, $n-1$ roads, all connected'' and quietly mean a
tree.
\end{insightbox}

\begin{theorem}{Unique Path}{c01-unique-path}
Between any two nodes of a tree there is exactly one simple path.
\end{theorem}
\begin{proofbox}
\emph{Existence:} the tree is connected, so some path exists between any two nodes.
\emph{Uniqueness:} suppose two different simple paths $P_1$ and $P_2$ joined $u$ to
$v$. Walk along $P_1$ from $u$ until the first node $w$ where $P_1$ and $P_2$ split,
then continue to where they rejoin at some node $z$. The two different $w$-to-$z$
pieces together form a closed loop, i.e.\ a cycle -- contradicting that the graph is a
tree. Hence the path is unique.
\end{proofbox}
\begin{insightbox}
Unique paths are why tree problems rarely need ``shortest'' in the statement: the one
path \emph{is} the shortest. They are also why we can talk about ``the'' path between
$u$ and $v$ without qualification, and why removing a single edge splits a tree into
exactly two pieces.
\end{insightbox}

\subsection{Storing a Tree in Memory}

Interview binary trees arrive as \code{TreeNode*} pointers (see \chref{c02}). Contest
trees arrive as a list of $n-1$ edges over nodes numbered $0..n-1$ or $1..n$; we store
them as an adjacency list and remember the parent so the walk never doubles back.

\begin{lstlisting}[caption={Build an adjacency list from edges and measure sizes with a parent-guard DFS}, label={lst:c01-build}]
#include <bits/stdc++.h>
using namespace std;

int n;
vector<vector<int>> adj;      // adj[v] = neighbours of v
vector<int> sz;               // sz[v] = size of subtree at v

// Root the tree at 0. p is the node we came from; never walk back to it.
void dfs(int v, int p) {
    sz[v] = 1;                       // count v itself
    for (int to : adj[v]) {
        if (to == p) continue;       // the parent guard: without this -> infinite loop
        dfs(to, v);
        sz[v] += sz[to];             // fold each child's subtree size upward
    }
}

int main() {
    n = 5;
    adj.assign(n, {});
    vector<pair<int,int>> edges = {{0,1}, {0,2}, {1,3}, {1,4}};
    for (auto [u, v] : edges) { adj[u].push_back(v); adj[v].push_back(u); }
    sz.assign(n, 0);
    dfs(0, -1);
    // edges = n - 1 = 4, and sz[0] == n, confirming the two theorems above.
}
\end{lstlisting}

\begin{complexitybox}
Building the adjacency list is $O(n)$ (each of the $n-1$ edges pushed twice). The DFS
visits every node once and every edge twice, so it is $O(n)$ in time and $O(n)$ memory
for the list plus the recursion stack. On a path-shaped tree the stack can be $n$ deep;
see \secref{c01-anatomy} for the iterative alternative.
\end{complexitybox}

\subsection{Worked Examples}

\begin{example}{A six-node tree}{c01-six}
Take the tree of \figref{c01-anatomy} and add a child $G$ under $C$. It now has $n=6$
nodes. Count the edges: $AB, AC, BD, BE, CF, CG$ -- that is $6$ nodes but only $5=6-1$
edges, matching \thmref{c01-edges-nodes}. The path from $D$ to $G$ is
$D \to B \to A \to C \to G$, length $4$, and it is the only such path
(\thmref{c01-unique-path}).
\end{example}

\begin{tracebox}
Running \lstname\ of \coderef{c01-build} on a path $0-1-2-3$ rooted at $0$ gives
$\size(3)=1$, $\size(2)=2$, $\size(1)=3$, $\size(0)=4$: each node's size is $1$ plus its
single child's size, so sizes count up from the far leaf.
\end{tracebox}

\subsection{Common Pitfalls}

\begin{warningbox}[Forgetting the parent guard]
A tree stored as an adjacency list has \emph{undirected} edges, so $B$ appears in
\code{adj[A]} and $A$ appears in \code{adj[B]}. A DFS that does not skip the node it came
from bounces $A\to B\to A\to B\ldots$ forever. The line \code{if (to == p) continue;} is
not an optimisation; it is what makes the traversal terminate.
\end{warningbox}

\begin{warningbox}[Connected but cyclic, or acyclic but split]
Both failures look tree-like. ``No cycle'' alone allows a forest; ``connected'' alone
allows a graph with a loop. Only together do they define a tree -- and only a tree has
the $n-1$ edge count. If your edge count is not $n-1$, you do not have a tree.
\end{warningbox}

\subsection{Solved Problems}

\solvedproblem{c01-same-tree}{Same Tree}{LeetCode 100}
  {https://leetcode.com/problems/same-tree/}{LC Easy}
  {tree definition, structural equality, recursion}

\begin{problemstatement}[Same Tree]
Given the roots \code{p} and \code{q} of two binary trees, decide whether they are
\emph{the same tree}: same shape \emph{and} same value at every node. Two empty trees
are the same; an empty tree and a non-empty one are not.
\end{problemstatement}

\solstep{Understanding the Problem}
``Same'' here is the precise structural equality the definition gives: the roots match in
value, the left subtrees are the same, and the right subtrees are the same. That is
already a recursive description, which is why a two-line recursion solves it and why a
flat comparison of node values is not enough -- shape matters.

\solstep{Key Observations}
\begin{observation}
Two trees are the same iff their roots carry the same value \emph{and} their left
subtrees are the same \emph{and} their right subtrees are the same.
\end{observation}
\begin{observation}
The base case is the \nil{} pointer: if either node is \nil, the trees match exactly
when \emph{both} are \nil{} at once.
\end{observation}

\solstep{Solution}
Recurse in lockstep on the two trees. At each pair of nodes, if both are null return
true; if exactly one is null return false; otherwise compare values and recurse on the
two children pairs.
\begin{figure}[H]
\centering
\begin{tikzpicture}[diagram]
  \begin{scope}[xshift=-1.3cm]
    \node[tnode] (a) at (0,0) {1};
    \node[tnode] (b) at (-1,-1.2) {2};
    \node[tnode] (c) at (1,-1.2) {3};
    \draw[tedge] (a) -- (b); \draw[tedge] (a) -- (c);
  \end{scope}
  \node[annot] at (0,-2) {tree p};
  \begin{scope}[xshift=2.3cm]
    \node[tnode] (a2) at (0,0) {1};
    \node[tnode] (b2) at (-1,-1.2) {2};
    \node[tnode] (c2) at (1,-1.2) {3};
    \draw[tedge] (a2) -- (b2); \draw[tedge] (a2) -- (c2);
  \end{scope}
  \node[annot] at (0,-2) {tree q};
\end{tikzpicture}
\caption{These two trees are the same: equal values and equal shape, so the answer is
\code{true}. Changing $q$'s right child $3$ to a left child would make the shapes
differ and the answer \code{false}.}\label{fig:c01-same-tree}
\end{figure}

\solstep{Proof of Correctness}
\begin{proofbox}[Induction on tree size]
By \thmref{c01-unique-path} a binary tree's structure is fully described by its root and
its two subtrees. Induct on $\size(p)+\size(q)$. Base: both empty -- the code returns
\code{p == q}, correct. Step: if exactly one is empty they differ and the code returns
false, correct. Otherwise both are non-empty: by the induction hypothesis the recursive
calls correctly decide whether the left subtrees match and whether the right subtrees
match; the code returns true exactly when the root values equal \emph{and} both recursive
answers are true -- precisely the definition of ``same tree''.
\end{proofbox}

\solstep{Implementation}
\begin{lstlisting}[caption={Same Tree (LeetCode 100)}, label={lst:c01-same-tree}]
#include <bits/stdc++.h>
using namespace std;

struct TreeNode {
    int val;
    TreeNode *left, *right;
    TreeNode(int x = 0) : val(x), left(nullptr), right(nullptr) {}
};

class Solution {
public:
    bool isSameTree(TreeNode* p, TreeNode* q) {
        if (p == nullptr || q == nullptr)   // both null -> equal; one null -> not
            return p == q;
        return p->val == q->val
            && isSameTree(p->left,  q->left)
            && isSameTree(p->right, q->right);
    }
};
\end{lstlisting}

\solstep{Complexity}
\begin{complexitybox}
Time $O(\min(n_p, n_q))$: we stop as soon as a mismatch or a null appears, visiting at
most the smaller tree. Memory $O(h)$ for the recursion stack, worst case $O(n)$ for a
skewed tree. Because the depth can be $n$, an iterative BFS-pair version (queue of node
pairs) avoids stack overflow on pathological inputs.
\end{complexitybox}

\solstep{Pitfalls and Variants}
\begin{warningbox}[Catching only value mismatches]
Comparing values but not shape passes a tree against its mirror image -- wrong. The
\code{p == q} null test is what enforces equal shape; do not replace it with two
separate \code{return true} branches that can drift apart.
\end{warningbox}
Variants: \emph{Symmetric Tree} (LC 101) compares a tree to its mirror; \emph{Subtree of
Another Tree} (LC 572) asks whether one tree appears inside another.

\subsection{Practice Problems}

\begin{practiceset}
  \practice{Invert Binary Tree}{LeetCode 226}
    {https://leetcode.com/problems/invert-binary-tree/}{LC Easy}
    {swap the two children at every node -- a one-line post-order recursion.}
  \practice{Univalued Binary Tree}{LeetCode 965}
    {https://leetcode.com/problems/univalued-binary-tree/}{LC Easy}
    {check every node equals the root; pass the expected value down in a DFS.}
  \practice{Leaf- Similar Trees}{LeetCode 865}
    {https://leetcode.com/problems/leaf-similar-trees/}{LC Easy}
    {collect leaf values left-to-right with a DFS, then compare the two lists.}
  \practice{Minimum Distance Between BST Nodes}{LeetCode 783}
    {https://leetcode.com/problems/minimum-distance-between-bst-nodes/}{LC Easy}
    {Stretch: an in-order walk of a BST is sorted -- compare neighbours. See \chref{c05}.}
\end{practiceset}
